% Prose rendering of PrimeTensor/Analysis/Examples.lean.
% Fragment: no \documentclass, no preamble, no document environment.

\section{Derived rules and positive-depth powers}
\label{Analysis:Examples:sec}

\leanfile{PrimeTensor/Analysis/Examples.lean}

\subsection*{What this file establishes}

Nothing in this file is new analysis, and the module header says so: it
develops the consequences that require no further continuity assumption.  Three
derivatives are computed --- of a ratio, of the reciprocal map, and of every
positive-depth power --- and all three are assembled from the product and
inversion rules of \texttt{PrimeTensor/Analysis/Rules.lean} and the two exact
examples of \texttt{PrimeTensor/Analysis/Derivative.lean}.

The substantive content is the power operation.  \texttt{PrimeTensor/Structure.lean}
rejected Lean's $\ctor{CommGroup}$ partly because it would equip the carrier
with a natural-power field defined at $0$; this file supplies the replacement,
a power indexed by $\name{Depth}$ with no zeroth case.  The power rule that
results, Theorem~\ref{Analysis:Examples:thm:power}, therefore has --- in the
docstring's words --- no zeroth case to prove or even state.

\subsection{Positive-depth powers}

\begin{definition}[\lean{MulReal.depthPow}]\label{Analysis:Examples:def:pow}
For $a : \name{MulReal}$, define $\name{depthPow}(a,-) : \name{Depth} \to
\name{MulReal}$ by
\[
  \name{depthPow}(a, \ctor{one}) = a,
  \qquad
  \name{depthPow}(a, \ctor{succ}\,d) = a \cdot \name{depthPow}(a, d),
\]
so that $\name{depthPow}(a,d) = a^{\size{d}}$.  Two $\ctor{simp}$ lemmas,
$\name{depthPow\_one}$ and $\name{depthPow\_succ}$, record the two equations,
both by $\ctor{rfl}$.
\end{definition}

\begin{leancode}
def depthPow (a : MulReal) : Depth → MulReal
  | .one => a
  | .succ d => a * depthPow a d
\end{leancode}

\begin{remark}[Two Depth-indexed iterations, doing different jobs]
\label{Analysis:Examples:rem:two-iterations}
This is the second operation in the development defined by recursion on a
depth, and it is worth setting beside the first.  $\name{scalePow}$ of
\texttt{PrimeTensor/Scale.lean} \emph{squares} at each successor and so reaches
exponent $2^{\size{d}-1}$ in $\size{d}$ steps; $\name{depthPow}$ multiplies by
one more copy and so reaches exponent $\size{d}$ in $\size{d}$ steps.  The
first climbs geometrically because it is measuring tolerances, which have to
become fine quickly; the second climbs linearly because it is an exponent,
which has to be able to take every positive value.  The same index type is
carrying two quite different scales, and only the linear one is a power in the
ordinary sense.
\end{remark}

\begin{designnote}[The zeroth power, promised absent and now absent]
\label{Analysis:Examples:note:no-zero}
\texttt{PrimeTensor/Structure.lean} declined $\ctor{CommGroup}$ on the ground
that $\ctor{Monoid}$ carries an $\ctor{npow}$ field defined at $0$, which would
reintroduce at the level of the carrier the zeroth index that
\texttt{PrimeTensor/Depth.lean} had excluded.  Declining a library operation is
only half a decision; this file completes it by defining the operation the
development actually wants.

The consequence is visible in Theorem~\ref{Analysis:Examples:thm:power}.  An
induction over $\mathbb{N}$ would have a base case $a^{0} = 1$ with derivative
$\name{trivial}$, and the statement would have to be checked there.  The
induction here starts at $\ctor{one}$, where the power is $a$ itself and the
derivative is $\name{identity}$ --- a case that carries content rather than
being a degenerate convention.  Nothing is lost: every exponent the additive
version could state, except $0$, is a depth.
\end{designnote}

\subsection{Derived tangent maps}

\begin{definition}[\lean{MulTangentMap.ratio}]\label{Analysis:Examples:def:tratio}
$\name{ratio}(D,E) \defeq \name{mul}(D, \name{inv}(E))$, with
$\name{ratio}(D,E)(h) = D(h)\,E(h)^{-1}$ by $\ctor{rfl}$.
\end{definition}

No verification is needed: the two operations of
\texttt{PrimeTensor/Analysis/Rules.lean} already produce tangent maps, so their
composite does too.  As the docstring notes, the ratio is built only from
product and inversion, introducing no third primitive on tangent maps --- the
same discipline \texttt{PrimeTensor/Structure.lean} applied in refusing
$\ctor{Div}$ on the carrier.

\begin{definition}[\lean{MulTangentMap.depthPow}]\label{Analysis:Examples:def:tpow}
\[
  \name{depthPow}(\ctor{one}) = \name{identity},
  \qquad
  \name{depthPow}(\ctor{succ}\,d) =
    \name{mul}(\name{identity},\, \name{depthPow}(d)).
\]
\end{definition}

\begin{leancode}
def depthPow : Depth → MulTangentMap
  | .one => identity
  | .succ d => mul identity (depthPow d)
\end{leancode}

\begin{lemma}[\lean{MulTangentMap.depthPow\_apply}]
\label{Analysis:Examples:lem:tpow-apply}
$\name{depthPow}(d)(h) = \name{depthPow}(h, d)$: the $d$-th power tangent map
sends a perturbation to its own $d$-th power.
\end{lemma}

\begin{leancode}
theorem depthPow_apply (d : Depth) (h : MulReal) :
    depthPow d h = MulReal.depthPow h d := by
  induction d with
  | one =>
      rfl
  | succ d ih =>
      change h * depthPow d h = h * MulReal.depthPow h d
      rw [ih]
\end{leancode}

\begin{proof}
By induction on $d$.  At $\ctor{one}$ both sides are $h$.  At
$\ctor{succ}\,d$ both sides unfold to $h$ times the respective value at $d$,
and the induction hypothesis identifies those.
\end{proof}

\begin{remark}[Stated but unused here]\label{Analysis:Examples:rem:unused}
Lemma~\ref{Analysis:Examples:lem:tpow-apply} is not cited by
Theorem~\ref{Analysis:Examples:thm:power}, whose proof unfolds both power
operations definitionally with $\ctor{change}$ and never needs the two to be
identified as terms.  The lemma is there to be quotable, and to record for the
reader that the two recursions --- one on the carrier, one on tangent maps ---
compute the same thing.
\end{remark}

\subsection{Three derivatives}

\begin{theorem}[\lean{hasMulDerivativeAt\_ratio}]
\label{Analysis:Examples:thm:ratio}
If $f$ has derivative $D_f$ at $x$ and $g$ has derivative $D_g$ at $x$, then
$y \mapsto \name{ratio}(f(y), g(y))$ has derivative
$\name{ratio}(D_f, D_g)$ at $x$.
\end{theorem}

\begin{proof}
Unfold both ratios to a product with an inverse and compose the two rules of
\texttt{PrimeTensor/Analysis/Rules.lean}.  The proof is one line because both
ratios were \emph{defined} as that composite, on the carrier and on tangent
maps alike; had either been a primitive, this would have been a third rule to
prove rather than a corollary.
\end{proof}

\begin{corollary}[\lean{hasMulDerivativeAt\_reciprocal}]
\label{Analysis:Examples:cor:reciprocal}
$y \mapsto y^{-1}$ has derivative $\name{inv}(\name{identity})$ at every point.
\end{corollary}

\begin{proof}
Apply the inversion rule to the identity derivative.
\end{proof}

\begin{theorem}[\lean{hasMulDerivativeAt\_depthPow}]
\label{Analysis:Examples:thm:power}
For every depth $d$ and every $x$, the map $y \mapsto \name{depthPow}(y,d)$ has
derivative $\name{depthPow}(d)$ at $x$.
\end{theorem}

\begin{proof}
By structural recursion on $d$.  At $\ctor{one}$ the map is the identity and
the claim is the identity derivative of
\texttt{PrimeTensor/Analysis/Derivative.lean}.  At $\ctor{succ}\,d$ the map is
$y \mapsto y \cdot \name{depthPow}(y,d)$ and the model is
$\name{mul}(\name{identity}, \name{depthPow}(d))$, so the product rule applies
to the identity derivative and the recursive call.
\end{proof}

\begin{designnote}[The derivative of a power does not depend on the point]
\label{Analysis:Examples:note:elasticity}
Theorem~\ref{Analysis:Examples:thm:power} gives one tangent map that works at
\emph{every} $x$, and by Lemma~\ref{Analysis:Examples:lem:tpow-apply} that map
is $h \mapsto h^{\size{d}}$.  The additive statement is
$(y^{n})' = n\,y^{\,n-1}$, which depends on $y$.  The discrepancy is not an
error on either side; it is the difference between two derivatives.

In the chart of \texttt{PrimeTensor/Analysis/Derivative.lean}, where both the
argument and the value are read logarithmically, the quantity being computed is
$\dfrac{d}{du}\Bigl[\log f\bigl(x e^{u}\bigr)\Bigr]_{u=0}
 = \dfrac{x f'(x)}{f(x)}$, the elasticity of $f$ at $x$.  For $f(y) = y^{n}$
that is the constant $n$, independent of $x$ --- which is exactly what the
theorem says, since $\log$ of $h \mapsto h^{n}$ is multiplication by $n$.

The whole family computed so far reads off in the same currency: elasticity $1$
for the identity, $0$ for a constant, $-1$ for the reciprocal, $n$ for the
$n$-th power, the sum of the elasticities for a product, and their difference
for a ratio.  Every rule proved in
\texttt{PrimeTensor/Analysis/Rules.lean} and in this file is the additivity of
elasticity, which is why none of them carries a Leibniz term and why none of
them mentions the point.  Nothing in this paragraph is formalised.
\end{designnote}

\begin{remark}[The negligibility machinery is still untouched]
\label{Analysis:Examples:rem:still-exact}
Every derivative known at the end of this file is obtained from
$\name{identity}$ and $\name{trivial}$ by products and inverses, so the
functions covered are exactly those of the form $y \mapsto c \cdot y^{n}$ with
$n$ an integer.  For such a function the response is
\[
  \frac{c\,(x h_k)^{n}}{c\,x^{n}} = h_k^{n},
\]
which is the model \emph{exactly}, with no error at all.  So every derivative
computed in the development so far routes through
$\name{firstOrder\_of\_exact}$, and $\name{NegligibleRelative}$ --- defined in
\texttt{PrimeTensor/Analysis/Derivative.lean} and equipped with closure
properties in \texttt{PrimeTensor/Analysis/Rules.lean} --- has still never been
applied to a nonzero error.  The first function needing it is not in this file.
\end{remark}

\begin{remark}[What is not proved]\label{Analysis:Examples:rem:missing}
$\name{depthPow}$ on the carrier is given no laws: neither
$a^{m}a^{n} = a^{m+n}$, nor $(ab)^{n} = a^{n}b^{n}$, nor
$(a^{m})^{n} = a^{mn}$ is stated, and $\name{Depth}$ carries no addition or
multiplication with which to state the first and third.  Negative powers exist
only as $(a^{n})^{-1}$, with no lemma identifying them with anything.  There is
still no chain rule and still no uniqueness of derivatives.
\end{remark}

\subsection*{Summary of the correspondence}

\begin{center}
\small
\begin{tabular}{@{}lll@{}}
\hline
Lean declaration & Kind & Here \\
\hline
\lean{MulReal.depthPow}                   & definition & Def.~\ref{Analysis:Examples:def:pow} \\
\lean{MulReal.depthPow\_one}              & simp thm   & Def.~\ref{Analysis:Examples:def:pow} \\
\lean{MulReal.depthPow\_succ}             & simp thm   & Def.~\ref{Analysis:Examples:def:pow} \\
\lean{MulTangentMap.ratio}                & definition & Def.~\ref{Analysis:Examples:def:tratio} \\
\lean{MulTangentMap.ratio\_apply}         & simp thm   & Def.~\ref{Analysis:Examples:def:tratio} \\
\lean{MulTangentMap.depthPow}             & definition & Def.~\ref{Analysis:Examples:def:tpow} \\
\lean{MulTangentMap.depthPow\_one\_apply} & simp thm   & Def.~\ref{Analysis:Examples:def:tpow} \\
\lean{MulTangentMap.depthPow\_apply}      & theorem    & Lem.~\ref{Analysis:Examples:lem:tpow-apply} \\
\lean{hasMulDerivativeAt\_ratio}          & theorem    & Thm.~\ref{Analysis:Examples:thm:ratio} \\
\lean{hasMulDerivativeAt\_reciprocal}     & theorem    & Cor.~\ref{Analysis:Examples:cor:reciprocal} \\
\lean{hasMulDerivativeAt\_depthPow}       & theorem    & Thm.~\ref{Analysis:Examples:thm:power} \\
\hline
\end{tabular}
\end{center}

\noindent
Every declaration of \texttt{PrimeTensor/Analysis/Examples.lean} appears above.
The file contains no \lean{sorry}, no axiom, and no unproved named proposition.
Design notes~\ref{Analysis:Examples:note:no-zero}
and~\ref{Analysis:Examples:note:elasticity} and
Remarks~\ref{Analysis:Examples:rem:two-iterations},
\ref{Analysis:Examples:rem:unused},
\ref{Analysis:Examples:rem:still-exact}
and~\ref{Analysis:Examples:rem:missing} are stated for the reader and are not
declarations of the file; the elasticity reading is a gloss and is not
formalised.
